\section{Las instituciones necesitan un Steering Wheel real}

Al final de la entrevista, las fallas de gobiernos y empresas se explican como una ruptura del feedback: los líderes tienen strategy y dashboard, pero no ven a tiempo lo que ocurre donde se ejecuta el trabajo; el personal de primera línea considera que los niveles superiores desconocen la realidad y traslada el trabajo real a Excel; los reportes agregados aparecen meses después y los errores se acumulan entre niveles jerárquicos. Sankar compara este estado con el volante falso del Jungle Cruise de Disneyland. Si el software quiere restituir la capacidad de las instituciones, debe conectar intent, execution, outcome y learning en un loop en tiempo real que pueda auditarse.

\begin{figure}[H]
\centering
\includegraphics[width=0.94\textwidth]{images/13-institutional-feedback.png}
\caption{Un steering wheel real requiere que objetivo, ejecución, resultado y aprendizaje formen un ciclo cerrado.}
\figsource{Redibujado a partir de los subtítulos 42:00--44:03, `images/13-institutional-feedback.png`.}
\end{figure}

\begin{importantbox}{Lectura de la figura: cuatro preguntas para verificar que el ciclo cerrado funciona}
Si el objetivo se convierte en restricciones ejecutables; si el estado de la ejecución es visible dentro del mismo sistema; si el outcome puede vincularse con una decision concreta; si el aprendizaje modifica de verdad la policy o el workflow de la siguiente ronda. Basta con que uno de estos pasos dependa de una spreadsheet fuera de línea, de reuniones verbales o de traslados manuales sin trazabilidad para que el volante se desconecte poco a poco de lo que ocurre en la operación.
\end{importantbox}

\begin{table}[H]
\centering
\begin{tabular}{L{0.19\textwidth}L{0.34\textwidth}L{0.37\textwidth}}
\toprule
Punto de ruptura & Síntoma visible & Capacidad que debe construirse \\
\midrule
Intent→ejecución & Los empleados no saben por qué cambian las prioridades & Policy, constraint y owner versionados \\
Ejecución→resultado & El dashboard solo muestra volumen de actividad, no resultados & Entity-level lineage y outcome metric \\
Resultado→aprendizaje & La revisión posterior no logra identificar qué decisión causó las consecuencias & Decision log, counterfactual y audit \\
Aprendizaje→Intent & Los incidentes del mismo tipo se repiten & Rule update, workflow test y rollout \\
Sistema formal→Excel & Las decisiones críticas salen de la plataforma & Modelado más rápido, interfaces explicables y cocreación con la primera línea \\
\bottomrule
\end{tabular}
\caption{Los puntos de ruptura del feedback institucional deben convertirse en capacidades de software y de organización observables.}
\end{table}

\begin{knowledgebox}{Nota de clase: no hay que pasar directamente de «falla institucional» a «demoler todo y empezar de cero»}
El ponente se opone a abandonar las instituciones solo porque la organización actual sea ineficiente. La tarea de la ingeniería de sistemas es restituir el sensing, la decision y el feedback, para que la policy pueda aprender de los resultados. Este juicio también necesita límites: reparar las instituciones no puede servir de justificación para ampliar un poder opaco, y el proceso de transformación debe conservar la supervisión, la división de poderes y la corrección de errores.
\end{knowledgebox}

\subsection{Resumen de la sección}

El valor último de la decision infrastructure no es un dashboard vistoso, sino que la organización cuente con un steering wheel real. Intent, execution, outcome y learning deben compartir entity, permisos y cadena de auditoría. El Excel escape, los reportes tardíos y los indicadores agregados que no pueden explicarse son señales de que el ciclo cerrado está distorsionado.
